\chapter{Boundary and interface conditions}
\label{ch:interface}

\section{Interface conditions on the fluid--solid boundary}
\label{sec:interface}
On the deformed interface $\Ga_{\mathrm{if}}$ the two subsystems exchange kinematics, tractions,
and heat. With the conventions of Section~\ref{sec:signs} the conditions are
\begin{align}
\bu\big(X+\bd(X)\big)&=0 &&\text{no-slip condition}, \label{eq:ifkin}\\
\bsig\cdot\bn &= \big(-p\,\bI+2\mu\mathbf D(\bu)\big)\cdot\bn &&\text{traction continuity}, \label{eq:iftraction}\\
T_f &= T_s &&\text{temperature continuity}, \label{eq:iftemp}\\
-k_f\nabla T_f\cdot\bn &= -k_s\nabla T_s\cdot\bn &&\text{heat-flux continuity}. \label{eq:ifflux}
\end{align}
The kinematic condition \eqref{eq:ifkin} is the steady form of the usual
$\bu=\partial\bd/\partial t$; the traction condition
\eqref{eq:iftraction} follows from balance of the interface traction vector in the absence of
interfacial stress, $\bsig^{(s)}\cdot\bn+\bsig^{(f)}\cdot\bn_f=0$ with $\bn_f=-\bn$, where
$\bsig^{(f)}=-p\bI+2\mu\mathbf D$ is the fluid stress. Evaluating both sides at the same point and
with the single normal $\bn$ yields \eqref{eq:iftraction}; the viscous part of the traction is
retained alongside the pressure, although at the baseline the pressure part dominates by about two
orders of magnitude in magnitude of its bending effect ($\tau_w=0.6$ Pa against
$\bar p=84$ Pa; Chapter~\ref{ch:baseline}). The conditions \eqref{eq:iftemp}--\eqref{eq:ifflux}
express perfect thermal contact (A8): continuity of temperature (no contact resistance) and
continuity of the normal heat flux, with the flux vector $\mathbf q=-k\nabla T$ evaluated with a
single normal, following the standard conjugate-heat-transfer formulation
\cite{carslaw1947conduction}.

\section{Exterior boundary conditions}
\label{sec:bcs}
Table~\ref{tab:bcs} specifies the remaining boundary data. The inlet profile is the fully developed
plane-Poiseuille parabola for a rigid entry of equal height, so that the mass flow per unit depth
is fixed at $\dot m'=\rho UH_0$ at both fixed-$U$ operation and (for the future fixed-$\Delta p$
variant) through the equivalent boundary data; the entrance length is a design input.

\begin{longtable}{p{1.0cm}p{2.6cm}p{7.2cm}p{3.9cm}}
\caption{Boundary-condition specification.}\label{tab:bcs}\\
\toprule
ID & Boundary & Condition & Remark \\
\midrule
\endfirsthead
\toprule
ID & Boundary & Condition & Remark \\
\midrule
\endhead
B1 & $x=-L_e$, fluid & $\bu=6U\,\tfrac{y}{H_0}\big(1-\tfrac{y}{H_0}\big)\hat{\mathbf x}$, $T_f=T_{\mathrm{in}}$ & fixed mass flow, rigid entry \\
B2 & $x=L+L_o$, fluid & traction-free outflow with datum $p=0$; $\partial_nT_f=0$ & exit-length/backflow sensitivity required (App.~B, C-9) \\
B3 & lower wall $y=0$ & $\bu=0$, $\partial_nT_f=0$ & rigid, adiabatic \\
B4 & rigid upper sections & $\bu=0$, $\partial_nT_f=0$ & rigid casing \\
B5 & $\Ga_{\mathrm{if}}$ & \eqref{eq:ifkin}--\eqref{eq:ifflux} & conjugate interface \\
B6 & outer face $\Ga_h^0$ & $-k_s\nabla T_s\cdot\bn_{\mathrm{out}}=q^{''}$ & baseline regime N \\
B7 & strip ends $X=0,L$ & $d_1=d_2=0$, zero applied moment; $\partial_nT_s=0$ & pinned, adiabatic (A9) \\
B8 & remaining solid faces & zero traction, zero flux & insulated casing \\
B9 & variant of B6 & $T_s=T_h$ on $\Ga_h^0$ & regime D; matched heat-input rule (A10) \\
\bottomrule
\end{longtable}

Two conditions deserve comment. The outflow condition B2 is the usual ``do-nothing'' natural
condition of finite-element practice for \eqref{eq:momentum}; it requires the exit length $L_o$ to
be sufficient that the flow has re-developed a nearly parallel state, verified a posteriori, and no
backflow occurs. The datum $p=0$ in B2, together with the casing reference of Assumption~A9,
fixes the gauge pressure carried by the deformable strip (Section~\ref{sec:deriv:pressure});
this datum is part of the problem specification and must be reported with every computed case. The outer-face condition B6 fixes the areal heat input and thereby the total heat
duty per unit depth, $Q'=q^{''}L$, which is the quantity balanced in the global enthalpy check of
Section~\ref{sec:quality}; regime D (B9) cannot share the heat-duty bookkeeping with regime N and
is compared only at matched $Q'$.

\section{Interface data flow in the coupled problem}
\label{sec:datamodel}
With the thermal load imposed on the solid side, the data flow of the coupled problem is: the fluid
subproblem receives its domain and the interface temperature from the solid; the solid subproblem
receives the interface traction (pressure-dominated) and the conjugate flux condition. Because the
solid problem is linear, the displacement decomposes exactly as
\begin{equation}
\bd=\bd_{\mathrm{th}}+\bd_p,
\label{eq:superposition}
\end{equation}
where $\bd_{\mathrm{th}}$ solves the thermoelastic problem with zero interface traction and
$\bd_p$ solves the isothermal problem under the fluid traction. This superposition is an algebraic
identity of the \emph{linear} solid model at fixed configuration and fixed loads; it is used as
the mechanism-attribution device between thermally driven and traction-driven deformation. Two
scope limits are recorded. First, across the coupled iterations the fluid state and the traction
distribution depend on $\bd$, so the decomposition is exact only at the converged state and for
the loads of that state. Second, at the von K\'arm\'an model level of
Section~\ref{sec:deriv:beamcolumn} the geometric-stiffness term couples the two load systems, so
the additive identity becomes approximate, with a relative error controlled by
$N_\theta^c/N_{\mathrm{cr}}$ ($\approx28\%$ in $w_{\max}$ at the baseline, from the
beam--column amplification); the implementation therefore attributes mechanisms by
\emph{separate load-case solutions} (thermal only, traction only, both) rather than by adding
fields, and verifies exact additivity on a detuned case at $N_\theta^c/N_{\mathrm{cr}}\le0.05$
(Chapter~\ref{ch:validation}).
